\scanpage{4}
\mspage{3}

\begin{center}- 57 -\end{center}
\margin{$\mathrm{n}^{\mathrm{o}}357$}

\[
f(x)\leq\sum_{i=1}^{n}a_i[\varphi(s_i)+\eta]g(s_i^{-1}x)
\]
et de même pour $f'$, et par suite :
\[
(f:g)+(f':g)\leq\sum_{i=1}^{n}a_i[\varphi(s_i)+\varphi'(s_i)+2\eta]\leq(1+2\eta)\sum_{i=1}^{n}a_i
\]
puisque $\varphi+\varphi'\leq1$. En appliquant la définition de $[F:g]$, puis (ii) et (vii), on en conclut :
\[
(f:g)+(f':g)\leq(1+2\eta)(F:g)\leq(1+2\eta)[(f+f':g)+\delta(h:g)]\ ;
\]
appliquant (v) à $f+f'$, $h$ et $g$, cela donne :
\[
(f:g)+(f':g)\leq(f+f':g)+2\eta(f+f':h)\cdot(h:g)+(1+2\eta)\delta\cdot(h:g).
\]
Si donc $f$, $f'$ et $h$ sont donnés comme dans (viii), il suffit de choisir $\delta$ et $\eta$ assez petits, puis de choisir $V$ comme il a été dit, pour qu'on ait l'inégalité formulée dans (viii).

Soit maintenant $\unsure{\mathfrak{F}|\mathcal{F}}$ un ultrafiltre sur $\mathcal{K}^{*}_{+}$, plus fin que le filtre engendré par la réunion des $\mathcal{K}^{*}_{+}(V)$ quand $V$ parcourt un système fondamental de voisinages compacts de $e$. D'autre part, fixons $f_0\in\mathcal{K}^{*}_{+}$ ; posons, pour $f\in\mathcal{K}^{*}_{+}$, $g\in\mathcal{K}^{*}_{+}$ :
\[
I_g(f)=\frac{(f:g)}{(f_0:g)},
\]
puis $I(f)=\lim_{\unsure{\mathfrak{F}|\mathcal{F}}}I_g(f)$. D'après (vi), $I(f)$ est dans l'intervalle $[1/(f_0:f),(f:f_0)]$. D'après (ii), on a $I(cf)=c.I(f)$ pour $c>0$. D'après (vii), on a $I(f+f')\leq I(f)+I(f')$. D'après (viii), on a $I(f+f')\geq I(f)+I(f')-\varepsilon I(h)$ quel que soit $\varepsilon\geq0$ si $h$ est $\geq1$ sur le support de $f+f'$ ; comme on peut toujours choisir $h$ ainsi dans $\mathcal{K}^{*}_{+}$, il s'ensuit que $I(f+f')=I(f)+I(f')$. Alors, d'après le Chap.II, § 2, n°1, prop.2, $I$ se prolonge (d'une manière et d'une seule) à une forme linéaire positive sur $\mathcal{K}$, c'est-à-dire qu'il y a une mesure positive $\alpha$ sur $G$ telle que $I(f)=\int f\,d\alpha$ pour $f\in\mathcal{K}^{*}_{+}$. D'après (i), on a $I(\gamma(\varepsilon_s)f)=I(f)$ quels que
